% Irodov, Part II section 2.4, original problems 140--145,
% printed pp.78--79, PDF pp.87--88.
\thermoentropyproblem{140}{78}
一摩尔范德瓦耳斯气体沿等温过程，
由体积 $V_1$ 变化到 $V_2$。
已知范德瓦耳斯参数 $a,b$，
求气体熵的变化 $\Delta S$。
\exsolution
一摩尔范德瓦耳斯状态方程为
\[
p=\frac{RT}{V-b}-\frac a{V^2}.
\]
对任一单相平衡物质，
Maxwell 关系给
\[
\left(\frac{\partial S}{\partial V}\right)_T
=\left(\frac{\partial p}{\partial T}\right)_V.
\]
在范德瓦耳斯气体中
$(\partial p/\partial T)_V=R/(V-b)$。
因此沿等温路径：
\[
\Delta S=\int_{V_1}^{V_2}
\frac R{V-b}\,\mathrm dV
=\boxed{R\ln\frac{V_2-b}{V_1-b}}.
\]
\excheck 引力修正参数 $a$ 在
$\partial p/\partial T$ 中被消去，
等温熵变只显式依赖排除体积参数 $b$；
条件是 $V_1,V_2>b$。

\thermoentropyproblem{141}{78}
一摩尔范德瓦耳斯气体从
初态 $(V_1,T_1)$ 到
末态 $(V_2,T_2)$，
其定容摩尔热容 $C_V$ 已知并视为常数。
求熵变。
\exsolution
将熵视为状态函数 $S(T,V)$，
其全微分为
\[
\mathrm dS=
\left(\frac{\partial S}{\partial T}\right)_V\mathrm dT+
\left(\frac{\partial S}{\partial V}\right)_T\mathrm dV.
\]
由定容热容定义
$(\partial S/\partial T)_V=C_V/T$，
由 Maxwell 关系与范德瓦耳斯状态方程
$(\partial S/\partial V)_T=R/(V-b)$，
因此
\[
\mathrm dS=\frac{C_V}{T}\,\mathrm dT+
\frac R{V-b}\,\mathrm dV.
\]
沿任意连接两状态的可逆路径积分：
\[
\boxed{\Delta S=C_V\ln\frac{T_2}{T_1}
+R\ln\frac{V_2-b}{V_1-b}}.
\]
\excheck 由于熵是状态函数，
实际过程即使不可逆，也可用连接初末平衡态的
假想可逆路径计算熵差。

\thermoentropyproblem{142}{78}
在很低温的区间中，
某晶体的总热容满足
$C(T)=aT^3$，其中 $a$ 为常数。
求晶体在此温区的熵
$S(T)$，并说明第三定律边界。
\exsolution
对可逆温度变化，
熵微分
$\mathrm dS=\delta Q_{\rm rev}/T
=C(T)\mathrm dT/T$。
代入 $C(T)=aT^3$：
\[
\mathrm dS=aT^2\,\mathrm dT.
\]
从参考温度 $T_0$ 积分：
\[
S(T)-S(T_0)
=\frac a3(T^3-T_0^3).
\]
若晶体基态唯一且满足第三定律
$S(0)=0$，可令 $T_0=0$，
于是
\[
\boxed{S(T)=\frac a3T^3}.
\]
\excheck 若存在残余熵，
一般表达式为
$S(T)=S(0)+aT^3/3$，
不能未经条件说明就强制设 $S(0)=0$。

\thermoentropyproblem{143}{78}
质量 $m=3.0\,\mathrm{kg}$ 的铝棒
从 $T_1=300\,\mathrm K$ 加热到
$T_2=600\,\mathrm K$。
在此温区其质量比热容为
$c(T)=a+bT$，其中
$a=0.77\,\mathrm{J/(g\,K)}$、
$b=0.46\,\mathrm{mJ/(g\,K^2)}$。
求棒的熵增。
\exsolution
对温度相关热容，
应积分 $\mathrm dS=mc(T)\mathrm dT/T$，
不能先求平均热容再乘 $\ln(T_2/T_1)$。
所以
\[
\Delta S=m\int_{T_1}^{T_2}
\frac{a+bT}{T}\,\mathrm dT
=m\left[a\ln\frac{T_2}{T_1}
+b(T_2-T_1)\right].
\]
取 $m=3000\,\mathrm g$、
$b=0.00046\,\mathrm{J/(g\,K^2)}$，
代入：
\[
\boxed{\Delta S=3000
[0.77\ln2+0.00046(300)]
\,\mathrm{J/K}
\simeq2.0\times10^3\,\mathrm{J/K}}.
\]
\excheck $b$ 的题给单位是
$\mathrm{mJ/(g\,K^2)}$，
须先乘 $10^{-3}$ 化为焦耳单位。

\thermoentropyproblem{144}{79}
某物质在一条过程中，
温度与熵的关系为
$T=aS^n$，其中 $a,n$ 均为常数。
求该过程的热容 $C(S)$，
并说明热容为负的条件。
\exsolution
热容由可逆热量定义：
\[
C=\frac{\delta Q_{\rm rev}}{\mathrm dT}
=T\frac{\mathrm dS}{\mathrm dT}.
\]
由 $T=aS^n$，
在 $S>0$、$n\ne0$ 的定义域：
\[
\frac{\mathrm dT}{\mathrm dS}
=anS^{n-1}=\frac{nT}{S}.
\]
所以
\[
\boxed{C(S)=\frac{T}{\mathrm dT/\mathrm dS}
=\frac{S}{n}}.
\]
对于正的绝对熵 $S>0$，
\[
\boxed{C<0\iff n<0}.
\]
\excheck 若 $n=0$，温度不随熵变化，
通常的有限热容表达式不适用；
热容为负的情形也必须结合体系
热力学稳定性和所处约束来理解。

\thermoentropyproblem{145}{79}
某物质沿一条\emph{热容恒为 $C$}
的多方路径变化。
当 $T=T_0$ 时熵为 $S_0$。
求温度作为熵的函数 $T(S)$，
并定性描绘 $C>0$ 与 $C<0$ 的图线。
\exsolution
由热容与熵关系
$C=T\,\mathrm dS/\mathrm dT$，
分离变量：
\[
\mathrm dS=C\,\frac{\mathrm dT}{T}.
\]
从 $(T_0,S_0)$ 积分到 $(T,S)$，
\[
S-S_0=C\ln(T/T_0).
\]
因此
\[
\boxed{T(S)=T_0
\exp\left(\frac{S-S_0}{C}\right)}.
\]
当 $C>0$，$T$ 随熵 $S$
单调指数增加；当 $C<0$，
$T$ 随熵单调指数下降，且二者均通过
点 $(S_0,T_0)$。
\begin{center}
\begin{tikzpicture}[scale=0.7,>=Stealth]
\draw[->] (0,0)--(4.8,0) node[right] {$S$};
\draw[->] (0,0)--(0,3.5) node[above] {$T$};
\draw[thick,domain=0.3:4.0,samples=45]
plot (\x,{0.40*exp(0.48*\x)});
\node at (4.1,3.2) {$C>0$};
\draw[thick,dashed,domain=0.3:4.2,samples=45]
plot (\x,{2.8*exp(-0.42*\x)});
\node at (3.65,0.38) {$C<0$};
\end{tikzpicture}
\end{center}
\excheck 熵沿路径增加不必然表示温度上升，
这一点由 $C$ 的符号决定；
图示只表达单调性，不按特定材料数值比例绘制。
